\subsection{Hamilton 同痕}\label{sec:10c}
下一个主题是\secref{8k}意义下延拓映射的高阶推广。这里只给出简要轮廓，略去技术细节。

在上面引入的 $\Ai$ 范畴 $\F(M)^{\mathrm{pr},\mathrm{free}}$ 中工作。设 $(\phi^s)$（$0\leq s\leq1$）是 $G=\operatorname{Aut}(M,\partial M)$ 中从 $\phi^0=\Id$ 到 $\phi^1=\phi$ 的同痕。按照\eref{8.23}的构造，可为每个 $(L,\gamma)$ 赋予一个上链
\[
T^0\in CF^{\mathrm{pr}}((L,\gamma),\phi(L,\gamma)).
\]
也就是说，取上半平面 $H$，配备由同痕 $\phi^s(L)$ 给出的移动 Lagrangian 边界条件，以及与给定 Floer 数据相容的扰动数据。注意，扰动数据的选择必然破坏 $H$ 的全纯自同构群。因此，将 $H$ 看作内部带有额外标点的曲面可能更自然，例如取 $z=i$，使曲面稳定。

沿着这个思路，更一般地考虑一个 $(d+1)$ 标点圆盘 $S$，它还带有额外的内部标点 $z$、一组条带状端，以及一组对象 $(L_k,\gamma_k)$（$k=0,\ldots,d$）。选择 $0=s_0\leq s_1\leq\cdots\leq s_{d+1}=1$，并为 $S$ 配备移动边界条件，使其在每个分支 $C_k$（$0\leq k\leq d$）上从 $\phi^{s_k}(L_k)$ 变化到 $\phi^{s_{k+1}}(L_k)$。然后在 $S$ 上选择扰动数据，使入向条带状端上的行为由对象对 $((L_0,\gamma_0),\phi(L_d,\gamma_d))$ 的 Floer 数据确定，出向端上的行为则分别由 $(\phi^{s_k}(L_{k-1},\gamma_{k-1}),\phi^{s_k}(L_k,\gamma_k))$ 的 Floer 数据确定。注意，由 $\F(M)^{\mathrm{pr},\mathrm{free}}$ 的定义，后者是将 $((L_{k-1},\gamma_{k-1}),(L_k,\gamma_k))$ 的 Floer 数据经 $\phi^{s_k}$ 诱导而得。

希望所有这些都光滑依赖于复结构及标点选择。因此，不应只考虑单个 $(S,z)$，而应考虑所有带额外一个内部标点的 $(d+1)$ 标点圆盘所成的模空间 $\mathcal R^{d+1,1}$，以及它上面的万有族。这个空间有 Deligne–Mumford 型紧化
\begin{equation}\label{eq:10.7}
\overline{\mathcal R}^{d+1,1}=\coprod_{(T,v_*)}\left(\mathcal R^{|v_*|,1}\times\prod_{v\ne v_*}\mathcal R^{|v|}\right),
\end{equation}
其各层以带有一个指定顶点 $v_*$ 的 $d$ 叶树为指标；$v_*$ 可以具有任意正度数，其他顶点则至少为三价。
\SourcePage{146}
条带状端、各 $\{s_k\}$ 及扰动数据的选择，必须满足与\eref{10.7}相关的适当一致性条件，这里不明确表述。随后，考虑对应的二元组 $(r\in\mathcal R^{d+1,1},u:\mathcal S_r^{d+1,1}\to M)$ 的模空间，其中 $u$ 是非齐次伪全纯映射。对该空间中满足适当渐近行为的点计数，定义映射
\[
\begin{aligned}
T^d:\ &CF^{\mathrm{pr}}((L_{d-1},\gamma_{d-1}),(L_d,\gamma_d))\otimes\cdots\\
&\cdots\otimes CF^{\mathrm{pr}}((L_0,\gamma_0),(L_1,\gamma_1))
\longrightarrow CF^{\mathrm{pr}}((L_0,\gamma_0),\phi(L_d,\gamma_d)).
\end{aligned}
\]
这里再次使用了以下事实：$CF^{\mathrm{pr}}(\phi^{s_k}(L_{k-1},\gamma_{k-1}),\phi^{s_k}(L_k,\gamma_k))$ 与 $CF^{\mathrm{pr}}((L_{k-1},\gamma_{k-1}),(L_k,\gamma_k))$ 可以典范地等同。各 $\{T^d\}$ 满足方程
\begin{equation}\label{eq:10.8}
\begin{aligned}
&\sum_{m,n}\mu^{d-m+1}\bigl(\phi_*^1(a_d),\ldots,\phi_*^1(a_{n+m+1}),
T^m(a_{n+m},\ldots,a_{n+1}),a_n,\ldots,a_1\bigr)\\
&\quad+\sum_{m,n}T^{d-m+1}\bigl(a_d,\ldots,a_{n+m+1},
\mu^m(a_{n+m},\ldots,a_{n+1}),a_n,\ldots,a_1\bigr)=0.
\end{aligned}
\end{equation}
\sr{fig:10.1}{图~10.1} 给出了两类求和项所对应的退化，并着重显示移动及固定的边界条件。由于 $\Ai$ 函子 $\phi_*$ 的高阶项为零，而且目前不需要考虑符号，\eref{10.8}恰好是 $T$ 成为从恒等函子到 $\phi_*$ 的自然变换的条件。像\secref{8k}末尾已经指出的，每个 $T^0$ 都是 $H\F(M)^{\mathrm{pr},\mathrm{free}}$ 中的同构，因此由\stref{引理}{1.6}，$T$ 在 $\Ai$ 函子的上同调范畴中为同构。将这个结果转移到与之拟等价的 $\Ai$ 范畴 $\F(M)^{\mathrm{pr}}$，得到：
\begin{proposition}\label{st:10.3}
设 $\phi\in\operatorname{Aut}(M,\partial M)$ 是与恒等映射同痕的自同构，则由 $\phi$ 诱导的 $\Ai$ 函子在 $H(\operatorname{fun}(\F(M)^{\mathrm{pr}},\allowbreak\F(M)^{\mathrm{pr}}))$ 中同构于恒等函子。
\end{proposition}
换言之，$\operatorname{Aut}(M,\partial M)$ 在 $\F(M)^{\mathrm{pr}}$ 上的协调作用，下降为 $\pi_0(\operatorname{Aut}(M,\partial M))$ 的弱作用。
\begin{note}\label{st:10.4}
考虑群 $\Gamma$ 在范畴 $A$ 上弱作用的一般情形。设 $\HH^0(A,A)^\times$ 是从 $\Id_A$ 到自身的函子同构所成的集合，以复合作为运算，它成为 Abel 群。将这一作用赋予协调结构的障碍，是一个群上同调类，更精确地说，是 $H^2(\Gamma;\HH^0(A,A)^\times)$ 中的元素。现在假设 $\Gamma=\pi_0(G)$、$G=\operatorname{Aut}(M,\partial M)$；为简便起见，取 $A=H\F(M)^{\mathrm{pr}}$ 为 Donaldson–Fukaya 范畴。此时，障碍类由 $H^2(\pi_0(G);\pi_1(G))$ 中一个典范类经同态 $\pi_1(G)\to\HH^0(A,A)^\times$ 诱导而来；这个典范类对任意拓扑群都存在，而同态利用\bcite{123}、\bcite{1}的方法定义。在 $\Ai$ 范畴层次，大概也存在相应陈述，
\SourcePage{147}
\SourceFigure{figure-10-1}{0.55869171}{10.1}
但我们尚未探究细节。无论如何，即使如此也仍然只是整个问题的一部分。与\stref{笔记}{10.1}一样，我们对 $\Ai$ 范畴上协调群作用的定义应当视为初步近似；可以期待，最终的概念将涉及 $G$ 的整个同伦型，而不仅仅是次数不超过 $1$ 的同伦群。
\end{note}

\subsection{接触型边界的情形}\label{sec:10d}
接续\secref{7b}，讨论具有接触型边界的正合辛流形。给定这样的流形 $(M,\varpi,\theta)$，可以将 $\mathcal J$ 的定义改为所有相容、且在边界附近为接触型的殆复结构所成的空间。这意味着，在所得 Floer 数据与扰动数据概念中，殆复结构族 $J(t)$ 或 $J(z)$ 在 $\partial M$ 附近不再保持不变。由\stref{引理}{7.3}，以及非齐次项在边界附近为零这一事实，方程\eref{8.8}、\eref{8.9}的解仍然留在 $M$ 的内部。以这种方式定义 $\F(M)^{\mathrm{pr}}$，会发现它本质上不依赖于任何殆复结构的选择；当然，这里的含义是不同选择组成协调系统。\znote{139}{这里用到接触型边界的凸性：$\partial M$ 附近可取 Liouville 结构，使 $\varpi=d(t\theta)$（$t\in(-\epsilon,0]$），于是 $d\theta$ 限制在 $J$-复直线上为正；对解作关于 $t$ 的最大值原理型论证，就排除了它触及边界。因此才可以在 $\partial M$ 附近自由改动 $J$ 而不改变模空间。\stref{引理}{7.3} 把这一论证表述为“解落在 $\operatorname{int}(M)$ 中”。[来源：原书 §(7b) 与\stref{引理}{7.3}，并见\eref{8.8}、\eref{8.9} 处的讨论。]}
\SourcePage{148}
\begin{proposition}\label{st:10.5}
设 $M$ 是具有接触型边界的正合辛流形，$N$ 是带角正合辛流形，$\iota:M\hookrightarrow\operatorname{int}(N)$ 是嵌入，满足 $\iota^*\varpi_N=\varpi_M$、$\iota^*\theta_N=\theta_M+d(\text{某个函数})$。则 $\F(M)^{\mathrm{pr}}$ 拟同构于 $\F(N)^{\mathrm{pr}}$ 的全 $\Ai$ 子范畴，后者由位于 $\iota(\operatorname{int}(M))$ 内的正合 Lagrangian 子流形组成。
\end{proposition}
\begin{proof}
可以给 $\theta_N$ 加上一个在 $\partial N$ 附近为零的函数的微分，而完全不影响 $\F(N)^{\mathrm{pr}}$。这样就将情形归结为 $\iota^*\theta_N=\theta_M$。此后，将 $\iota$ 简单视为包含映射。在 $N$ 上固定殆复结构 $J$，使其在 $\partial N$ 附近与 $I_N$ 相同，且 $J|_M$ 在 $\partial M$ 附近为接触型。

定义所讨论的 $\F(N)^{\mathrm{pr}}$ 的子范畴时，可以选择所有殆复结构在 $N\setminus\operatorname{int}(M)$ 的一个邻域中等于 $J$，并选择 Hamilton 函数与扰动数据，使它们在这个集合的一个邻域中为零。于是可应用\stref{引理}{7.5}，说明所有非齐次伪全纯多边形都留在 $\operatorname{int}(M)$ 中。在此情形，两个范畴实际上相等。对于一般的数据选择，拟同构结论由前面关于良定义性的讨论得到。
\end{proof}
\begin{corollary}\label{st:10.6}
设 $M,N$ 是具有接触型边界的正合辛流形，$\widehat M,\widehat N$ 是像\eref{7.1}那样接上锥得到的非紧致流形。假设存在正合辛同构 $\widehat\phi:\widehat M\to\widehat N$，则范畴 $\F(M)^{\mathrm{pr}}$、$\F(N)^{\mathrm{pr}}$ 拟等价。\znote{140}{要点是：$\F(M)^{\mathrm{pr}}$ 实际只依赖接触型边界流形的“完成”$\widehat M=M\cup_{\partial M}([0,\infty)\times\partial M)$，而 $\widehat\phi$ 是完成之间的正合辛同构。证明方式是把\stref{命题}{10.5}用两次：$M$ 嵌入 $\widehat M$、$N$ 嵌入 $\widehat N$，而 $\widehat\phi$ 把 $\iota_M(\operatorname{int}M)$ 送到 $\iota_N(\operatorname{int}N)$ 的紧致子集内，于是 $\F(M)^{\mathrm{pr}}$ 与 $\F(N)^{\mathrm{pr}}$ 都与 $\widehat M$ 上的同一个全子范畴拟同构。[来源：原书 §(10d)，完成与\eref{7.1} 的定义见 §(7b)。]}
\end{corollary}
\begin{proof}
考虑完备化的有限部分
\[
\widehat M^\sigma=M\cup_{\partial M}([0,\sigma]\times\partial M),\qquad
\widehat N^\tau=N\cup_{\partial N}([0,\tau]\times\partial N),
\]
选取 $0\ll\sigma\ll\tau$，使 $N\subset\widehat\phi(\widehat M^\sigma)\subset\widehat N^\tau$。首先，$M$ 的 Fukaya 范畴与 $\widehat M^\sigma$ 的拟等价；类似地，$N$ 的与 $\widehat N^\tau$ 的拟等价。这很容易看出：利用负 Liouville 流，两个流形可在辛形式及其原始一形式作共形缩放后彼此识别。若相应地缩放全部 Floer 数据及扰动数据，两者的 Fukaya 范畴实际上相同。

现在应用\stref{命题}{10.5}，它说明 $\F(\widehat M^\sigma)^{\mathrm{pr}}$ 拟同构于 $\F(\widehat N^\tau)^{\mathrm{pr}}$ 的某个全子范畴。该子范畴显然包含所有位于 $N$ 内的正合 Lagrangian 子流形，而在同痕意义下，这已经包含一切，因此得到所需拟等价。
\end{proof}
由\stref{引理}{7.2}，得到：
\begin{corollary}\label{st:10.7}
在具有接触型边界的正合辛流形这一类中，$\F(M)^{\mathrm{pr}}$ 在 $M$ 的变形下不变。
\end{corollary}
\SourcePage{149}
\begin{note}\label{st:10.8}
稍作变化，也可以使用完备化上自然的一类殆复结构 $\widehat J$，直接定义范畴 $\F(\widehat M)^{\mathrm{pr}}$。由\stref{引理}{7.4}可知，$\F(M)^{\mathrm{pr}}$ 与 $\F(\widehat M)^{\mathrm{pr}}$ 拟等价。这种做法的一个优点，是得到在无穷远处为接触型的正合辛自同构所成的群 $\operatorname{Aut}(\widehat M)$ 在 $\F(M)^{\mathrm{pr}}$ 上的协调作用。
\end{note}
\begin{corollary}\label{st:10.9}
设 $N$ 是带角正合辛流形。像\stref{引理}{7.6}那样，在充分小的 $\delta$ 处将角磨圆，得到具有接触型边界的正合辛流形 $M\subset\operatorname{int}(N)$。则 $\F(M)^{\mathrm{pr}}$ 与 $\F(N)^{\mathrm{pr}}$ 拟等价。
\end{corollary}
\begin{proof}
这立即由\stref{命题}{10.5}及每个正合 Lagrangian 子流形都能移入 $M$ 内部这一事实推出。
\end{proof}
\begin{corollary}\label{st:10.10}
在拟等价意义下，带角正合辛流形的 Fukaya 范畴在变形下不变。
\end{corollary}
\begin{proof}
取 $N$ 及一个变形 $(\varpi_N^s,\theta_N^s,I_N^s)$。考察\stref{引理}{7.6}的证明，可见存在 $M\subset\operatorname{int}(N)$，使得对所有 $s$，$(\varpi_N^s,\theta_N^s)$ 在 $M$ 上的限制都给出具有接触型边界的正合辛流形结构。\stref{推论}{10.7}说明 $\F(M)^{\mathrm{pr}}$ 不依赖于 $s$，而\stref{推论}{10.9}使这一结果可以转移到 $\F(N)^{\mathrm{pr}}$。
\end{proof}
\begin{corollary}\label{st:10.11}
设 $\phi:N\to\widetilde N$ 是带角正合辛流形的同构。则存在从 $\F(N)^{\mathrm{pr}}$ 到 $\F(\widetilde N)^{\mathrm{pr}}$ 的拟等价，将每个正合 Lagrangian 子流形 $L$ 映到 $\phi(L)$。
\end{corollary}
在给出证明之前，应先澄清一个可能产生混淆的地方。若同构与殆复结构相容，即边界附近有 $\phi_*(I_N)=I_{\widetilde N}$，则\secref{10a}中的形式论证就足以证明结论。事实上，这些论证给出稍强的陈述，因为所得拟等价在相关函子范畴中自动在同构意义下典范。一般情形的困难，恰恰来自\secref{7a}所定义的同构概念不考虑殆复结构。
\begin{proof}
将角磨圆，得到具有接触型边界的正合辛流形 $M\subset\operatorname{int}(N)$。限制 $\phi|_M$ 显然是正合辛嵌入。取\stref{推论}{10.9}中拟等价的拟逆，再应用\stref{命题}{10.5}，得到 $\Ai$ 函子
\begin{equation}\label{eq:10.9}
\F(N)^{\mathrm{pr}}\longrightarrow\F(M)^{\mathrm{pr}}\longrightarrow\F(\widetilde N)^{\mathrm{pr}},
\end{equation}
\SourcePage{150}
它是上同调全忠实的。将 $\delta$ 取得充分小，可使 $\widetilde N$ 的任意正合 Lagrangian 子流形都能通过负 Liouville 流移入 $\phi(\operatorname{int}(M))$，于是\eref{10.9}就是拟等价。按定义，它把 $L$ 映到 $\phi(\widetilde L)$，其中 $\widetilde L$ 与 $L$ 同痕。但因为 $\phi(L)$、$\phi(\widetilde L)$ 是 $H\F(\widetilde N)^{\mathrm{pr}}$ 中同构的对象，纯代数论证便可修改此函子，使 $L$ 确实映到 $\phi(L)$。
\end{proof}
